\chapter{Dirac Equation}

\section*{Introduction}

Dirac wanted an equation that was both quantum and relativistic. The Schr\"odinger equation was quantum but not relativistic; the Klein--Gordon equation was relativistic but gave nonsensical results for electrons. Dirac found the unique equation satisfying both requirements---but it had ``extra'' solutions describing particles with positive energy but opposite charge. These were positrons, discovered by Carl Anderson in 1932. The equation also revealed that the electron has intrinsic angular momentum (spin) not due to spinning but to the structure of spacetime itself.


The Dirac equation for a free particle of mass $m$:
\[
(i\gamma^\mu \partial_\mu - m)\psi = 0
\]
where $\gamma^\mu$ ($\mu = 0,1,2,3$) are $4\times 4$ gamma matrices satisfying $\{\gamma^\mu, \gamma^\nu\} = 2g^{\mu\nu}I$, and $\psi$ is a 4-component spinor.

\section*{Fundamental Equations Used}

The derivation starts from the relativistic energy--momentum relation $E^2 = p^2 c^2 + m^2 c^4$ and the requirement for a first-order Lorentz-covariant quantum equation.

\section*{Assumptions}

\textbf{Assumptions:} The particle is spin-1/2. Spacetime is flat. The particle is free; for interactions, replace $\partial_\mu$ with $D_\mu = \partial_\mu + ieA_\mu$.


\textbf{Limitations:} Does not include gravity (coupling to curved spacetime is possible but not renormalizable). Does not describe particle creation/annihilation (need QFT). Breaks down near the Planck scale.


\section*{Mathematical Derivation}

\textbf{Step 1: The problem with the Klein--Gordon equation.}

The relativistic energy--momentum relation is $E^2 = p^2c^2 + m^2c^4$. The naive quantization $E \to i\hbar\partial_t$, $\mathbf{p} \to -i\hbar\nabla$ gives the Klein--Gordon equation:
\begin{equation}
\frac{1}{c^2}\frac{\partial^2 \phi}{\partial t^2} - \nabla^2 \phi + \frac{m^2c^2}{\hbar^2}\phi = 0
\end{equation}
This is second-order in time, leading to negative probability densities---physically unacceptable for a single-particle equation.

\textbf{Step 2: Dirac's requirement---first order in time.}

Dirac required the equation to be first-order in both space and time (like the Schr\"odinger equation), yet Lorentz covariant. The general form must be:
\begin{equation}
i\hbar\frac{\partial \psi}{\partial t} = \left(c\,\boldsymbol{\alpha} \cdot \hat{\mathbf{p}} + \beta m c^2\right)\psi
\end{equation}
where $\hat{\mathbf{p}} = -i\hbar\nabla$ and $\boldsymbol{\alpha} = (\alpha_1, \alpha_2, \alpha_3)$ and $\beta$ are to be determined.

\textbf{Step 3: Recover the relativistic dispersion relation.}

Square this equation (acting on a free-particle solution $e^{-i(Et - \mathbf{p}\cdot\mathbf{r})/\hbar}$):
\begin{equation}
E^2 = c^2\mathbf{p}^2\,\boldsymbol{\alpha}^2 + m^2c^4\,\beta^2 + mc^3\mathbf{p}\cdot(\boldsymbol{\alpha}\beta + \beta\boldsymbol{\alpha})
\end{equation}
For this to equal $E^2 = c^2\mathbf{p}^2 + m^2c^4$, we require:
\begin{equation}
\alpha_i^2 = 1,\quad \beta^2 = 1,\quad \alpha_i\alpha_j + \alpha_j\alpha_i = 2\delta_{ij},\quad \alpha_i\beta + \beta\alpha_i = 0
\end{equation}

\textbf{Step 4: Identify the Clifford algebra.}

These are the anticommutation relations of the Clifford algebra:
\begin{equation}
\{\gamma^\mu, \gamma^\nu\} = 2g^{\mu\nu}I_4
\end{equation}
where $\gamma^0 = \beta$, $\gamma^i = \beta\alpha_i$, and $g^{\mu\nu} = \text{diag}(+1,-1,-1,-1)$. The minimal representation requires $4\times 4$ matrices (Pauli's theorem), giving a 4-component spinor $\psi$.

\textbf{Step 5: Choose a representation and write the covariant form.}

One common representation uses:
\begin{equation}
\gamma^0 = \begin{pmatrix} I_2 & 0 \\ 0 & -I_2 \end{pmatrix}, \quad \gamma^i = \begin{pmatrix} 0 & \sigma_i \\ -\sigma_i & 0 \end{pmatrix}
\end{equation}
where $\sigma_i$ are the Pauli matrices. The covariant form is:
\begin{equation}
i\hbar\gamma^\mu\partial_\mu\psi - mc\,\psi = 0
\end{equation}
In natural units ($\hbar = c = 1$):
\begin{equation}
\boxed{(i\gamma^\mu\partial_\mu - m)\psi = 0}
\end{equation}

\textbf{Step 6: Identify the spinor structure.}

The 4-component spinor decomposes as:
\begin{equation}
\psi = \begin{pmatrix} \psi_A \\ \psi_B \end{pmatrix}
\end{equation}
where $\psi_A$ has two components describing positive-energy states (electrons) and $\psi_B$ describes negative-energy states (positrons). In the non-relativistic limit, $\psi_B \to 0$ and the equation reduces to the Pauli equation with $g = 2$ for the electron---spin emerges from the Dirac equation without being added by hand.

\section*{Characteristics of the Equation}

\begin{itemize}
\item \textbf{Spin emerges naturally:} Unlike Schr\"odinger (where spin is added by hand), the Dirac equation produces spin from Lorentz covariance.
\item \textbf{Antimatter prediction:} Negative-energy solutions interpreted as positrons, confirmed by Anderson in 1932.
\item \textbf{Fine structure:} Predicts hydrogen energy levels including correct $2S_{1/2}$--$2P_{1/2}$ degeneracy.
\item \textbf{Magnetic moment:} Predicts electron $g$-factor $g = 2$ (QED corrections give $g = 2.0023\ldots$).
\item \textbf{Zitterbewegung:} Rapid oscillatory motion at the Compton wavelength from positive/negative energy interference.
\end{itemize}
\section*{Solution Methods}

\begin{enumerate}
\item \textbf{Plane waves:} $\psi = u(p) e^{-ip \cdot x}$, with $(\gamma^\mu p_\mu - m) u = 0$. Two positive-energy (electron) and two negative-energy (positron) solutions.
\item \textbf{Foldy--Wouthuysen transformation:} Unitary transformation separating positive/negative energies, reducing to approximate Schr\"odinger form with relativistic corrections.
\item \textbf{Numerical methods:} Lattice QCD solves the equation on a spacetime grid to compute hadron properties.
\item \textbf{Bound states:} Hydrogen solution uses spherical symmetry, yielding exact energy levels.
\end{enumerate}
\section*{Verifications}

\begin{itemize}
\item \textbf{Positron discovery:} Anderson's 1932 discovery confirmed Dirac's prediction.
\item \textbf{Hydrogen fine structure:} Measured levels agree with Dirac to better than $10^{-6}$.
\item \textbf{Electron $g$-factor:} $g_e/2 = 1.001\,159\,652\,181\,643(764)$, measured to 12 digits, matching QED built on Dirac.
\item \textbf{Graphene:} Measured linear dispersion and universal optical absorption ($\pi\alpha \approx 2.3\%$) match massless Dirac predictions.
\end{itemize}
\section*{Applications}

\begin{itemize}
\item \textbf{Atomic physics:} Precise predictions of hydrogen energy levels, fine structure, hyperfine structure.
\item \textbf{Particle physics:} The Standard Model is built on the Dirac equation for quarks and leptons.
\item \textbf{PET imaging:} Medical imaging uses positrons (predicted by Dirac) from radioactive tracers.
\item \textbf{Graphene:} Electrons behave as massless Dirac fermions, producing extraordinary electronic properties.
\item \textbf{Topological insulators:} Surface states described by the Dirac equation lead to novel spintronic applications.
\end{itemize}
\section*{What Richard Feynman Would Have Asked}

\subsection*{1. Plain-English Translation}
\textit{``What is the Dirac equation saying, without any jargon?''}

It says: ``An electron is a wave that obeys both quantum mechanics and Einstein's relativity, and the mathematical requirements force it to have spin and an antimatter twin.'' The equation is not optional---there is only one way to write a relativistic wave equation for a spin-1/2 particle, and this is it. Spin and antimatter are not added by hand; they are consequences of mathematical consistency.

The Dirac equation is like a lock that only one key fits. The ``lock'' is the requirement of consistency with both quantum mechanics and special relativity. The ``key'' is the Dirac equation itself, with its four-component spinor, gamma matrices, and prediction of both electron and positron states.

\subsection*{2. Localized Mechanics}
\textit{``What is the electron doing at each point in space, according to the Dirac equation?''}

The Dirac spinor $\psi$ has four components at each point: two for electron spin up/down, two for positron spin up/down. The probability of finding the electron with a given spin is $|\psi_i|^2$ for the appropriate component. Spin is not ``located'' at a point---it is a property of the wave function at that point.

He would press you on spin and relativity: in the non-relativistic limit, the Dirac equation reduces to the Pauli equation with spin added. But in the full relativistic equation, spin emerges from the gamma matrix structure---the same matrices ensuring Lorentz covariance. Spin is not classical; it is a relativistic quantum property with no classical analogue.

\subsection*{3. Testing Extreme Limits}
\textit{``What happens at extremely high energies or in extremely strong fields?''}

At very high energies ($E \gg mc^2$), the mass term becomes negligible and the equation reduces to the Weyl equation for massless spin-1/2 particles. This is the particle physics regime where electrons, quarks, and neutrinos are effectively massless.

In extremely strong electric fields ($E \sim 10^{18}$\,V/m), the vacuum becomes unstable and electron--positron pairs are spontaneously created (Schwinger effect). At the Planck scale ($10^{19}$\,GeV), gravity becomes important and the Dirac equation must be generalized to curved spacetime via the spin connection.

\subsection*{4. Dimensionality and Geometry}
\textit{``Does the Dirac equation care about the dimensionality of space?''}

Yes, profoundly. In 3+1 dimensions, $4\times 4$ gamma matrices and 4-component spinors are needed. In 2+1 dimensions (graphene), $2\times 2$ matrices suffice---physics differs (no chiral symmetry breaking, different quantum Hall effect). In 1+1 dimensions, the equation simplifies further and can be solved exactly.

On curved spacetime, the equation is modified by the spin connection, coupling the spinor to geometry. On a lattice (lattice QCD), careful discretization avoids ``doubling''---spurious fermion species from naive discretization.

\subsection*{5. Cross-Disciplinary Connection}
\textit{``Where else does the Dirac equation appear outside of particle physics?''}

In condensed matter, the Dirac equation appears wherever electrons behave as relativistic fermions. In graphene, electrons near Dirac points obey the 2D massless Dirac equation, producing extraordinary electronic properties. In topological insulators, surface states are described by the Dirac equation, leading to protected conducting states robust against disorder. In the quantum Hall effect, edge states are chiral Dirac fermions.

The mathematical structure---a first-order differential equation for a multi-component wave function with Clifford algebra symmetry---appears in any system where two quantum states are coupled by a gradient. It describes neutrino oscillations, Majorana fermions, and certain models of cosmic strings.


%=============================================================
% Chapter 21: Electromagnetic Wave Equation
%=============================================================
\section*{FAQ}

\textbf{Q1: What exactly is spin, if the electron is a point particle?}

A: Spin is not the electron ``spinning''---a point particle cannot spin. It is an intrinsic angular momentum, a fundamental quantum property arising from the requirement that the Dirac equation be Lorentz covariant. The electron simply has spin-1/2 as a property of its existence, just as it has charge $-e$.

\textbf{Q2: Why are the gamma matrices $4\times 4$ and not $2\times 2$?}

A: The minimal Lorentz group representation for spin-1/2 requires four components: two for the electron (spin up/down) and two for the positron (spin up/down). A $2\times 2$ representation (Pauli matrices) works non-relativistically but cannot incorporate both particle and antiparticle states.

\textbf{Q3: Can the Dirac equation describe quarks?}

A: Yes, it applies to all spin-1/2 fermions. The difference is that quarks carry color charge, so $eA_\mu$ is replaced by the QCD gauge field $A_\mu^a T^a$. The basic Dirac structure is the same.
\section*{Important Bibliography}

\begin{enumerate}
\item Dirac, P.A.M., ``The Quantum Theory of the Electron,'' \textit{Proceedings of the Royal Society A} 117, 610--624 (1928).
\item Bjorken, J.D. and Drell, S.D., \textit{Relativistic Quantum Mechanics} (McGraw-Hill, 1964), Chapter 1.
\item Peskin, M.E. and Schroeder, D.V., \textit{An Introduction to Quantum Field Theory} (Addison-Wesley, 1995), Chapter 3.
\item Greiner, W., \textit{Relativistic Quantum Mechanics: Wave Equations}, 3rd ed. (Springer, 2000), Chapter 7.
\end{enumerate}

